Da una breve historia de Kali Linux.

El desarrollo de Kali Linux comenzó con un proyecto llamado \textbf{Whoppix} (\textit{White Hat Knoppix}), que era una distribución Linux en formato Live CD diseñada para realizar auditorías de seguridad y pruebas de penetración. Esta distribución estaba basada en Knoppix, un sistema operativo ejecutado desde un disco óptico sin necesidad de instalación, e incluía herramientas para la seguridad informática, análisis de redes y exploits.

Posteriormente, el siguiente proyecto pasó a estar basado en Slax, llamado \textbf{WHAX} (\textit{WhiteHat Slax}). El proyecto surgió de una fusión con el proyecto Auditor Security Collection. A partir de esta alianza surgió BackTrack Linux, basado en Slackware. 

En 2012 comenzó el desarrollo de Kali Linux a partir del deseo de reemplazar BackTrack Linux, que era un proyecto mantenido manualmente, con el objetivo de ser un derivado de Debian con toda la infraestructura necesaria y con técnicas de empaquetado mejoradas. Se decidió que Kali Linux estuviera construida sobre Debian por la calidad y la estabilidad que ofrece. La primera versión de Kali Linux se lanzó en marzo de 2013 basada en Debian 7~\cite{debstory}.
